\section{Preliminaries}\label{section:prelim}
We recap some preliminaries on $\two$-categories, simplicial objects, Segal properties, and monads and adjunctions. 

\subsection{Higher category theory}\label{subsection:catthy}
We collect together some higher category theory we use throughout the paper. The properties we need are largely independent of the choice of a model of $\one$ or $\two$-category. We assume to have fixed nested Grothendieck universes which we refer to as small, large and very large.

All $\one$-categories appearing as the coefficients of our perverse schobers will be stable presentable large categories. All functors between $\one$-categories are assumed to be colimit-preserving, or equivalently admit (a possibly non-colimit-preserving) right adjoint. We denote by $\sfSt_k$ the $\two$-category of stable large presentable categories with colimit-preserving functors and natural transformations.

We denote by $\wCat_{\two}$ the $\one$-category of large $\two$-categories. As pointed out in \cite{stefanichPresentable$inftyN$categories2020}, $\sfSt_k$ is only large, roughly because presentable categories are determined by a small amount of data. Thus $\sfSt_k\in \wCat_{\two}$ defines an object.

Given a $\two$-category $\sfA\in\wCat_{\two}$, and objects $\calA_1,\calA_2\in\sfA$, we denote by $\sfA(\calA_1,\calA_2)$ the mapping category. Given also $\sfB\in\wCat_{\two}$, we say that a pair of functors $\opname{L}:\sfA\tofrom \sfB:\opname{R}$ are adjoint if there exist functorial equivalences of mapping categories,
\[\sfA(\calA,\opname{R}(\calA))\simeq \sfB(\opname{L}(\calA),\calB).\]

Given a functor $\opname{I}:\sfA\to \sfB$, we say that
\begin{itemize}
    \item $\opname{I}$ is 2-fully faithful if $\opname{I}$ induces an equivalence on all mapping categories,
    \item $\opname{I}$ is 2-faithful if $\opname{I}$ induces a fully faithful functor on all mapping categories,
    \item $\opname{I}$ is a locally full inclusion if $\opname{I}$ is 2-faithful and induces a monomorphism on equivalence classes of objects.
\end{itemize}

Whenever $\frakd$ is a small $\two$-category, there is a $\two$-category $\sfFun(\frakd,\sfSt_k)\in \wCat_{\two}$ whose objects are functors $\frakd\to \sfSt_k$. Whenever $f:\frakd_1\to \frakd_2$ is a functor between $\two$-categories, there is a restriction functor,
\[f^*:\sfFun(\frakd_2,\sfSt_k)\to \sfFun(\frakd_1,\sfSt_k).\]

All $\two$-categories of schobers will be constructed out of these diagram categories $\sfFun(\frakd,\sfSt_k)\in \wCat_{\two}$ by taking locally full sub-$\two$-categories. All functors will be constructed using either Theorem~\ref{thm:laxfibrations} or Proposition~\ref{prop:functorconstruction}, which we discuss below.

\begin{rmk}
    We expect that all $\two$-categories of schobers are ``$k$-linear stable presentable $\two$-categories'' and that all functors respect this structure. A potential candidate for such theory is laid out in \cite{stefanichPresentable$inftyN$categories2020}. In this paper, we will content ourselves with working inside $\wCat_{\two}$.
\end{rmk}

\begin{thm}[Abell\'an--Haugseng--Martini, \cite{abellanFreeFibrationsLax2026}*{Theorem 5.6.5}]\label{thm:laxfibrations}
    The functor $f^*$ admits both adjoints $f_!\adjto f^*\adjto f_*$ which each have explicit descriptions on objects.
\end{thm}

Our other method of constructing functors is as follows.

\begin{prop}\label{prop:functorconstruction}
    Let $\sfC,\sfD,\sfS\in\wCat_{\two}$, and let $\opname{F}:\sfC\to\sfS$, and $\opname{G}:\sfD\to\sfS$ be functors, so that we have lifts $\sfC,\sfD\in (\wCat_{\two})_{/\sfS}$. Assume that $\opname{G}$ is 2-faithful. Then, the mapping space
    \[\Map_{(\wCat_{\two})_{/\sfS}}(\sfC,\sfD)\]
    is discrete. Moreover, any functor $\sfC\to\sfD$ over $\sfS$ is determined by the map it induces on the sets of equivalence classes of objects.
\end{prop}

We note that the proof of the proposition will also give us a criterion under which a given map on equivalence classes of objects comes from a functor. We will need the following standard facts.

\begin{lem}\label{lem:ftoff}
    If a functor $\opname{G}:\sfD\to\sfS$ is 2-faithful, then the diagonal functor $\opname{Diag}:\sfD\to\sfD\times_{\sfS}\sfD$ is 2-fully faithful.
\end{lem}

\begin{proof}
    Let $\calD_1,\calD_2\in\sfD$. We have an identification
    \[(\sfD\times_{\sfS}\sfD)(\opname{Diag}\calD_1,\opname{Diag}\calD_2)\simeq \sfD(\calD_1,\calD_2)\times_{\sfS(\opname{G}\calD_1,\opname{G}\calD_2)} \sfD(\calD_1,\calD_2),\]
    and the functor $\sfD(\calD_1,\calD_2)\to (\sfD\times_{\sfS}\sfD)(\opname{Diag}\calD_1,\opname{Diag}\calD_2)$ is the diagonal functor. The diagonal of a fully faithful functor of $\one$-categories is an equivalence, so we are done.
\end{proof}

Lemma~\ref{lem:ftoff} allows us to reduce statements about 2-faithful functors to those about 2-fully faithful ones which are much easier to think about. In the following, we recall that for $k\geq -1$ a map $f$ between spaces is said to be $k$-truncated if the diagonal of $f$ is $(k-1)$-truncated, and that a map $f$ is said to be $(-2)$-truncated if it is an equivalence. So a $(-1)$-truncated map is an inclusion of connected components, and a $0$-truncated map is one which is relatively discrete.

\begin{lem}\label{lem:discrete1}
    Let $\frakt\in\wCat_{\two}$ be any $\two$-category. Then, the post-composition map on mapping categories
    \[\opname{G}\circ-:\Map(\frakt,\sfD)\to \Map(\frakt,\sfS)\]
    is $0$-truncated.
\end{lem}

\begin{proof}
    This is equivalent to showing that the diagonal map of $\opname{G}\circ-$ is $(-1)$-truncated. This diagonal map is identified with $\opname{Diag}\circ-$, and by Lemma~\ref{lem:ftoff}, $\opname{Diag}$ is 2-fully faithful. We are done by observing that $\opname{Diag}\circ-$ is $(-1)$-truncated. In words, a functor $\frakt\to \sfD\times_{\sfS}\sfD$, if it factors through the full sub-$\two$-category $\sfD$, does so uniquely.
\end{proof}

\begin{lem}\label{lem:discrete2}
    Let $\frakt'\to\frakt$ be a functor of $\two$-categories which is essentially surjective. Then, the natural map
    \[\Map(\frakt,\sfD)\to \Map(\frakt',\sfD)\times_{\Map(\frakt',\sfS)}\Map(\frakt,\sfS)\]
    is $(-1)$-truncated.
\end{lem}

\begin{proof}
    This is equivalent to showing that the diagonal is $(-2)$-truncated, i.e. an equivalence. This is identified with the natural map
    \[\Map(\frakt,\sfD)\to \Map(\frakt',\sfD)\times_{\Map(\frakt',\sfD\times_{\sfS}\sfD)}\Map(\frakt,\sfD\times_{\sfS}\sfD).\]
    By Lemma~\ref{lem:ftoff}, $\opname{Diag}$ is 2-fully faithful, so that this is indeed an equivalence. In words, a functor $\frakt\to \sfD\times_{\sfS}\sfD$ factors through $\sfD$ after restricting to $\frakt'$, if and only if it factors through $\sfD$.
\end{proof}

\begin{proof}[Proof of Proposition~\ref{prop:functorconstruction}]
    By Yoneda lemma, to give a functor $\sfC\to \sfD$ over $\sfS$ is equivalent to giving for each $\frakt\in\wCat_{\two}$ a commuting diagram of spaces,
    % https://q.uiver.app/#q=WzAsMyxbMCwwLCJcXE1hcChcXGZyYWt0LFxcc2ZDKSJdLFsxLDAsIlxcTWFwKFxcZnJha3QsXFxzZkQpIl0sWzEsMSwiXFxNYXAoXFxmcmFrdCxcXHNmUykiXSxbMSwyXSxbMCwyXSxbMCwxLCIiLDAseyJzdHlsZSI6eyJib2R5Ijp7Im5hbWUiOiJkYXNoZWQifX19XSxbNCwxLCJcXHNpbSIsMSx7InNob3J0ZW4iOnsic291cmNlIjoyMH0sInN0eWxlIjp7ImJvZHkiOnsibmFtZSI6ImRhc2hlZCJ9fX1dXQ==
    \[\begin{tikzcd}[ampersand replacement=\&]
        {\Map(\frakt,\sfC)} \& {\Map(\frakt,\sfD)} \\
        \& {\Map(\frakt,\sfS)}
        \arrow[dashed, from=1-1, to=1-2]
        %\arrow["\opname{F}\circ-"{name=0, anchor=center, inner sep=0}, from=1-1, to=2-2]
        \arrow["\opname{F}\circ-"', from=1-1, to=2-2]
        \arrow["\opname{G}\circ-", from=1-2, to=2-2]
        %\arrow["\sim"{description}, between={0.2}{1}, Rightarrow, dashed, from=0, to=1-2]
    \end{tikzcd}\]
    which is functorial in $\frakt$.
    
    By Lemma~\ref{lem:discrete1}, the arrow $\opname{G}\circ-$ is relatively discrete, and hence the space of lifts of $\opname{F}\circ-$ is discrete.

    Suppose we exhibit a lift for the trivial $\two$-category $\frakt=\mathfrak{pt}$. This determines the lift for $\frakt$ any discrete space, as well as the functoriality for any map between discrete spaces.
    
    We show that in fact, the lift is uniquely determined for all $\frakt\in \wCat_{\two}$, if it exists. To see this, let $\frakt'\to\frakt$ be any essentially surjective functor from any discrete space $\frakt'$. Functoriality gives us a commutative diagram,
    % https://q.uiver.app/#q=WzAsNixbMCwxLCJcXE1hcChcXGZyYWt0LFxcc2ZDKSJdLFsxLDEsIlxcTWFwKFxcZnJha3QsXFxzZkQpIl0sWzEsMiwiXFxNYXAoXFxmcmFrdCxcXHNmUykiXSxbMSwwLCJcXE1hcChcXGZyYWt0JyxcXHNmQykiXSxbMiwwLCJcXE1hcChcXGZyYWt0JyxcXHNmRCkiXSxbMiwxLCJcXE1hcChcXGZyYWt0JyxcXHNmUykiXSxbMSwyXSxbMCwyXSxbMCwxLCIiLDAseyJzdHlsZSI6eyJib2R5Ijp7Im5hbWUiOiJkYXNoZWQifX19XSxbNCw1XSxbMyw1XSxbMyw0XSxbMCwzXSxbMSw0XSxbMiw1XV0=
    \[\begin{tikzcd}[ampersand replacement=\&]
        \& {\Map(\frakt',\sfC)} \& {\Map(\frakt',\sfD)} \\
        {\Map(\frakt,\sfC)} \& {\Map(\frakt,\sfD)} \& {\Map(\frakt',\sfS)} \\
        \& {\Map(\frakt,\sfS)}
        \arrow[from=1-2, to=1-3]
        \arrow[from=1-2, to=2-3]
        \arrow[from=1-3, to=2-3]
        \arrow[from=2-1, to=1-2]
        \arrow[dashed, from=2-1, to=2-2]
        \arrow[from=2-1, to=3-2]
        \arrow[from=2-2, to=1-3]
        \arrow[from=2-2, to=3-2]
        \arrow[from=3-2, to=2-3]
    \end{tikzcd}\]
    The desired lift denoted by the dotted arrow is uniquely determined if it exists, by Lemma~\ref{lem:discrete2} applied to the square on the right of this diagram.
    
    Suppose now that lifts exist for all $\frakt\in \wCat_{\two}$. We claim that the functoriality of these lifts is automatic. Let $\frakt_1\to \frakt_2$ be any map of $\two$-categories. Choose any map $\frakt_1'\to \frakt_2'$ of discrete spaces together with a commutative diagram as follows, such that the vertical maps are essentially surjective.
    \[\begin{tikzcd}[ampersand replacement=\&]
        {\frakt_1'} \& {\frakt_2'} \\
        {\frakt_1} \& {\frakt_2}
        \arrow[from=1-1, to=1-2]
        \arrow[from=1-1, to=2-1]
        \arrow[from=1-2, to=2-2]
        \arrow[from=2-1, to=2-2]
    \end{tikzcd}\]
    A similar argument to before produces the desired commuting diagram
    % https://q.uiver.app/#q=WzAsNCxbMCwwLCJcXE1hcChcXGZyYWt0XzEsXFxzZkMpIl0sWzEsMCwiXFxNYXAoXFxmcmFrdF8xLFxcc2ZEKSJdLFsxLDEsIlxcTWFwKFxcZnJha3RfMixcXHNmRCkiXSxbMCwxLCJcXE1hcChcXGZyYWt0XzIsXFxzZkMpIl0sWzAsMV0sWzMsMF0sWzIsMV0sWzMsMl1d
    \[\begin{tikzcd}[ampersand replacement=\&]
        {\Map(\frakt_1,\sfC)} \& {\Map(\frakt_1,\sfD)} \\
        {\Map(\frakt_2,\sfC)} \& {\Map(\frakt_2,\sfD).}
        \arrow[from=1-1, to=1-2]
        \arrow[from=2-1, to=1-1]
        \arrow[from=2-1, to=2-2]
        \arrow[from=2-2, to=1-2]
    \end{tikzcd}\]
    Higher functoriality is similarly automatic. We deduce that $\Map_{(\wCat_{\two})_{/\sfS}}(\sfC,\sfD)$ embeds as a subset of the set of lifts for $\frakt=\mathfrak{pt}$, proving the proposition.
\end{proof}

\subsection{Simplicial objects}\label{subsection:simplicialintro}
We denote by $\Delta$ the simplex category, whose objects are $[n] := \{0,\dots,n\}$ for $n\geq0$ and whose morphisms are non-decreasing functions $f:[m]\to[n]$. We consider the following standard set of generating morphisms.
\begin{itemize}
    \item For each $n\geq1$ let $\delta_{n,j}:[n-1]\to[n]$ be the injective map skipping $j$.
    \item For each $n\geq0$ let $\sigma_{n,j}:[n+1]\to[n]$ be the surjective map hitting $j$ twice.
\end{itemize}

We say that a morphism $f:[m]\to [n]$ is active if $f(0)=0$ and $f(m)=n$. We say that a morphism $f:[m]\to [n]$ is inert if $f$ is injective and its image is convex under the partial order. These two classes of morphisms are each closed under composition and so define subcategories $\Delta_{\act}$ and $\Delta_{\inert}$ respectively. Note that in $\Delta_{\act}$, $[1]$ is an initial object. These two classes of morphisms in fact form a factorization system of $\Delta$, that is, every morphism $f:[m]\to [n]$ uniquely factors as $i\circ a$ where $a$ is active and $i$ is inert.

Let $\calC$ be a $\one$-category with fibered products. A simplicial object $C(-)$ in $\calC$ is a functor
\[C:\Delta^{\op}\to \calC.\]
We often write $C(n) := C([n])$, as well as
\begin{align*}
    d_{n,j} &:= C(\delta_{n,j}): C(n)\to C(n-1)\\
    s_{n,j} &:= C(\sigma_{n,j}): C(n)\to C(n+1).
\end{align*}

A simplicial object $C(-)$ is said to be 1-Segal if for every pushout square of inert maps
\[\begin{tikzcd}[ampersand replacement=\&]
	{[n]}\pullback \& {[n-m]} \\
	{[m]} \& {[0]}
	\arrow[from=1-2, to=1-1]
	\arrow[from=2-1, to=1-1]
	\arrow[from=2-2, to=1-2]
	\arrow[from=2-2, to=2-1]
\end{tikzcd}\]
the resulting square of objects in $\calC$ is a pullback square. By definition, this upgrades $C(1)$ to a monoidal object in the slice category $\calC_{/C(0)}$.

Many of our simplicial objects will be 2-Segal, a notion introduced by Dyckerhoff--Kapranov \cite{dyckerhoffHigherSegalSpaces2019}, and independently under the name decomposition spaces by G\'alvez-Carrillo--Kock--Tonks \cite{galvez-carrilloDecompositionSpacesIncidence2018}.

A simplicial object $C(-)$ is said to be 2-Segal if for every pushout square of maps with $a$ active and $i$ inert (which implies $a'$ active and $i'$ inert)
\[\begin{tikzcd}[ampersand replacement=\&]
	{[n]} \& {[n+1-m]} \\
	{[m]} \& {[1]}
	\arrow["{i'}"', from=1-2, to=1-1]
	\arrow["{a'}", from=2-1, to=1-1]
	\arrow["a", from=2-2, to=1-2]
	\arrow["i"', from=2-2, to=2-1]
\end{tikzcd}\]
the resulting square of objects in $\calC$ is a pullback square.

\begin{eg}\label{eg:waldhausen}
    The prototypical example of such an object arises from the Waldhausen S-construction, of which the following is a simple instance. Let $\calC$ be the category of sets, and let $C(n)$ denote the set of filtered vector spaces $0=V_0\subset V_1\subset\cdots \subset V_n$. We first define the face maps
    \[d_{n,i}:C(n)\to C(n-1).\]
    For $i=0$, this outputs $V_1/V_1\subset V_2/V_1\subset\cdots\subset V_n/V_1$, for $i=n$, this outputs $0=V_0\subset V_1\cdots \subset V_{n-1}$, and for $0<i<n$, this outputs the filtered vector space where we remove $V_i$ from the filtration. We define the degeneracy maps
    \[s_{n,i}:C(n)\to C(n+1).\]
    For each $0\leq i\leq n$, this outputs the result of repeating $V_i$ twice. The 2-Segal condition reduces to some elementary linear algebra.
\end{eg}

We describe a mechanism we will frequently use in order to produce simplicial $\two$-categories and maps between simplicial $\two$-categories from an existing one.

\begin{lem}\label{lem:descendingsimplicialobjects}
    Let $\sfC'(-):\Delta^{\op}\to \wCat_{\two}$ be a simplicial $\two$-category, and suppose for each $n$, we have a sub-$\two$-category $\opname{I}_n:\sfC(n)\inclto \sfC'(n)$ admitting a right adjoint $\opname{I}_n^R$.

    For each arrow $a:[m]\to[n]$ in $\Delta$, define
    \[\sfC(a) := \opname{I}_m^R\circ\sfC'(a)\circ \opname{I}_n,\]
    so that we have a lax-commuting square as follows.
    \[\begin{tikzcd}[ampersand replacement=\&]
        {\sfC'(n)} \& {\sfC'(m)} \\
        {\sfC(n)} \& {\sfC(m)}
        \arrow["{\sfC'(a)}", from=1-1, to=1-2]
        \arrow["{\opname{I}_n^R}"', from=1-1, to=2-1]
        \arrow["{\opname{I}_n^R}", from=1-2, to=2-2]
        \arrow[between={0}{0.8}, Rightarrow, from=2-1, to=1-2]
        \arrow["{\sfC(a)}"', from=2-1, to=2-2]
    \end{tikzcd}\]
    If these actually commute for all arrows, or equivalently a generating set of arrows, then $\sfC(-)$ defines a simplicial $\two$-category,
    \[\sfC(-):\Delta^{\op}\to\wCat_{\two},\]
    equipped with a functor, $\sfC'(-)\to \sfC(-)$.
\end{lem}

\begin{proof}
    We allow ourselves to enhance the target and consider the $\two$-category of $\two$-categories, $\widehat{\sfCat}_{\two}$. As defined, we have a lax-functor $\sfC:\Delta^{\op}\to \widehat{\sfCat}_{\two}$. The conditions ensure that this lax-functor factors through $\wCat_{\two}$. Indeed, let $[m]\xto{a}[n]\xto{b}[p]$ be morphisms with composition $c$. Then the 2-morphism $\sfC(a)\circ\sfC(b) \Rightarrow \sfC(c)$ can be written as
    \begin{align*}
        \sfC(a)\circ\sfC(b) &\simeq \sfC(a)\circ\sfC(b)\circ\opname{I}_p^R\circ \opname{I}_p\\
        &\simeq \sfC(a)\circ\opname{I}_n^R\circ\sfC'(b)\circ\opname{I}_p\\
        &\simeq \opname{I}_m^R\circ\sfC'(a)\circ\sfC'(b)\circ\opname{I}_p\\
        &\simeq \sfC(c).
    \end{align*}
    Since $\Delta^{\op}$ is a 1-category, $\sfC$ factors through the maximal $\one$-category,
    \[\wCat_{\two}\inclto \widehat{\sfCat}_{\two}.\]
\end{proof}

A similar mechanism will allow us to produce simplicial maps from existing ones.

\begin{lem}\label{lem:descendingsimplicialmaps}
    Let $\opname{F}'(-):\sfC'(-)\to \sfD'(-)$ be a map between two simplicial $\two$-categories, and suppose that for each $n$, we have sub-$\two$-categories
    \[\opname{I}_n:\sfC(n)\inclto \sfC'(n), \opname{J}_n:\sfD(n)\inclto \sfD'(n)\]
    admitting right adjoints $\opname{I}_n^R,\opname{J}_n^R$.

    Suppose that the conditions of Lemma~\ref{lem:descendingsimplicialobjects} are satisfied for both $\sfC'(-)$ and $\sfD'(-)$ so that we have simplicial $\two$-categories $\sfC(-),\sfD(-)$ as well as maps $\sfC'(-)\to \sfC(-)$ and $\sfD'(-)\to \sfD(-)$.
    
    For each arrow $n$, define
    \[\opname{F}(n):= \opname{J}_n^R\circ\opname{F}'(n)\circ \opname{I}_n: \sfC(n)\to\sfD(n)\]
    so that we have a lax-commuting square as follows.
    % https://q.uiver.app/#q=WzAsNCxbMCwwLCJcXHNmQycobikiXSxbMSwwLCJcXHNmRCcobikiXSxbMSwxLCJcXHNmRChuKSJdLFswLDEsIlxcc2ZDKG4pIl0sWzEsMl0sWzAsM10sWzAsMSwiXFxvcG5hbWV7Rn0nKG4pIl0sWzMsMiwiXFxvcG5hbWV7Rn0obikiLDJdLFszLDEsIiIsMCx7InNob3J0ZW4iOnsidGFyZ2V0IjoyMH0sImxldmVsIjoyfV1d
    \[\begin{tikzcd}[ampersand replacement=\&]
        {\sfC'(n)} \& {\sfD'(n)} \\
        {\sfC(n)} \& {\sfD(n)}
        \arrow["{\opname{F}'(n)}", from=1-1, to=1-2]
        \arrow[from=1-1, to=2-1]
        \arrow[from=1-2, to=2-2]
        \arrow[between={0}{0.8}, Rightarrow, from=2-1, to=1-2]
        \arrow["{\opname{F}(n)}"', from=2-1, to=2-2]
    \end{tikzcd}\]
    If these actually commute for all $n$, then $\opname{F}(-)$ defines a map of simplicial $\two$-categories,
    \[\opname{F}(-):\sfC(-)\to\sfD(-),\]
    and there is a commuting diagram of simplicial $\two$-categories, as follows.
    \[\begin{tikzcd}[ampersand replacement=\&]
        {\sfC'(-)} \& {\sfD'(-)} \\
        {\sfC(-)} \& {\sfD(-)}
        \arrow["{\opname{F}'(-)}", from=1-1, to=1-2]
        \arrow[from=1-1, to=2-1]
        \arrow[from=1-2, to=2-2]
        \arrow["{\opname{F}(-)}"', from=2-1, to=2-2]
    \end{tikzcd}\]
\end{lem}

\begin{proof}
    The proof is similar to that of Lemma~\ref{lem:descendingsimplicialobjects}.
\end{proof}

\subsection{Monads and adjunctions}\label{subsection:mndadj}
We review the theory of monads and adjunctions following \cites{riehlHomotopyCoherentAdjunctions2015,haugsengLaxTransformationsAdjunctions2021}.

We will define the $\two$-category $\sfAdj$ of adjunctions in $\sfSt_k$ as a locally full subcategory of $\sfFun(\frakn(1),\sfSt_k)$, where $\frakn(1)$ is the small $\two$-category $1\from \infty$. This locally full subcategory consists of objects
\[\calA_{1}\xto{F}\calA_{\infty}\]
where $F$ is a left adjoint functor, and morphisms
% https://q.uiver.app/#q=WzAsNCxbMCwwLCJcXGNhbEFfMSJdLFsxLDAsIlxcY2FsQV9cXGluZnR5Il0sWzAsMSwiXFxjYWxBXzEnIl0sWzEsMSwiXFxjYWxBX1xcaW5mdHknIl0sWzAsMSwiRiJdLFsyLDMsIkYnIl0sWzAsMiwiXFxvcG5hbWV7R31fMSIsMl0sWzEsMywiXFxvcG5hbWV7R31fXFxpbmZ0eSJdXQ==
\[\begin{tikzcd}[ampersand replacement=\&]
	{\calA_1} \& {\calA_\infty} \\
	{\calA_1'} \& {\calA_\infty'}
	\arrow["F", from=1-1, to=1-2]
	\arrow["{\opname{G}_1}"', from=1-1, to=2-1]
	\arrow["{\opname{G}_\infty}", from=1-2, to=2-2]
	\arrow["{F'}", from=2-1, to=2-2]
\end{tikzcd}\]
where the commuting square is adjointable, that is, the natural transformation $\opname{G}_1F^R\to (F')^R\opname{G}_\infty$ is an equivalence.

It is shown in \cite{riehlHomotopyCoherentAdjunctions2015} that there is a small $\two$-category $\frakadj$ called the walking adjunction, such that $\sfFun(\frakadj,\sfSt_k)$ also parametrizes adjunctions in $\sfSt_k$. It follows from \cite{haugsengLaxTransformationsAdjunctions2021} that these two $\two$-categories in fact agree i.e. the functor obtained by restriction to the left adjoint $f:\frakn(1)\to \frakadj$ defines an equivalence,
\[f^*:\sfFun(\frakadj,\sfSt_k)\isomto \sfAdj.\]

This second description of $\sfAdj$ is useful for studying monads. Let $i:\frakmnd\inclto \frakadj$ be the full sub-$\two$-category of $\frakadj$ generated by the source of the left adjoint. Then,
\[\sfMnd:=\sfFun(\frakmnd,\sfSt_k)\]
parametrizes pairs $(\Phi,M)$ where $\Phi$ is an object of $\sfSt_k$ and $M$ is a monad.

\begin{rmk}\label{rmk:monadstrictness}
    We note that morphisms $(\Phi_1,M_1)\to (\Phi_2,M_2)$ are given by a functor $F_{21}:\Phi_1\to \Phi_2$ together with an identification $F_{21}M_1\simeq M_2 F_{21}$. In particular the fiber over $\Phi\in\sfSt_k$ of $\sfMnd$ has no non-invertible morphisms. To see such morphisms, we must allow morphisms in $\sfFun$ to be given by lax natural transformations, as studied in \cite{haugsengLaxTransformationsAdjunctions2021}. This distinction becomes important in \S\ref{subsection:suspension}.
\end{rmk}

There is a pullback functor $i^*:\sfAdj\to \sfMnd$ which on objects sends the adjunction $S:\Phi\tofrom\Psi:R$ to $(\Phi,RS)$. By Theorem~\ref{thm:laxfibrations}, this admits a left and right adjoint, 
\[i_!,i_*:\sfMnd\to \sfAdj.\]
These are the stable, presentable analogs of the Kleisli and Eilenberg--Moore adjunctions respectively. Note that
\[i^*i_*\isomto \Id\]
so that these each embed $\sfMnd$ as a full sub-$\two$-category of $\sfAdj$.

We recall that classically, the Kleisli adjunction is constructed by restricting to the essential image of the left adjoint in the Eilenberg--Moore adjunction. A key observation in \cite{christSphericalMonadicAdjunctions2023} is that the essential image of a functor of stable categories need not be stable. By passing to a stable closure, we obtain a good notion of Kleisli adjunction. It is shown in \cite{gammagePerverseSchobers3d2023}*{Proposition B.4} that this notion agrees with the left adjoint $i_!$.

In our presentable setting, we must also account for the fact that the essential image of a continuous functor of presentable categories need not be presentable. We formulate an analogous recognition principle in our setting.

\begin{lem}\label{lem:kleislirecognition}
    An adjunction $S:\Phi\tofrom\Psi:R$ is in the sub-$\two$-category $i_!(\sfMnd)$ if and only if the smallest stable, presentable subcategory which contains $\opname{Im}(S)$ is equal to $\Psi$.
\end{lem}

\begin{proof}
    The proof is identical to the proof of \cite{gammagePerverseSchobers3d2023}*{Proposition B.4}.
\end{proof}

From $i^*i_*\simeq \Id$, we obtain a natural transformation
\[i_!\to i_*,\]
which is an analog of the classical comparison between the Kleisli and Eilenberg--Moore adjunctions. In our presentable setting, this natural transformation is an equivalence.

\begin{prop}\label{prop:KLisEM}
    The natural transformation $i_!\to i_*$ is an equivalence.
\end{prop}

\begin{proof}
    Let $(\Phi,M)\in\sfMnd$ and denote $i_!(\Phi,M)\to i_*(\Phi,M)$ as follows.
    % https://q.uiver.app/#q=WzAsNCxbMCwwLCJcXFBoaSJdLFsxLDAsIlxcb3BuYW1le0tMfShGKSJdLFswLDEsIlxcUGhpIl0sWzEsMSwiXFxNb2QoRikiXSxbMCwxLCJTX3tcXG9wbmFtZXtLTH19IiwwLHsib2Zmc2V0IjotMX1dLFswLDIsIlxcaWRfe1xcUGhpfSIsMl0sWzIsMywiU197XFxvcG5hbWV7RU19fSIsMCx7Im9mZnNldCI6LTF9XSxbMSwzLCJJIl0sWzEsMCwiUl97XFxvcG5hbWV7S0x9fSIsMCx7Im9mZnNldCI6LTF9XSxbMywyLCJSX3tcXG9wbmFtZXtFTX19IiwwLHsib2Zmc2V0IjotMX1dXQ==
    \[\begin{tikzcd}
        \Phi & {\opname{KL}(M)} \\
        \Phi & {\Mod(M)}
        \arrow["{S_{\opname{KL}}}", shift left, from=1-1, to=1-2]
        \arrow["{\id_{\Phi}}"', from=1-1, to=2-1]
        \arrow["{R_{\opname{KL}}}", shift left, from=1-2, to=1-1]
        \arrow["I", from=1-2, to=2-2]
        \arrow["{S_{\opname{EM}}}", shift left, from=2-1, to=2-2]
        \arrow["{R_{\opname{EM}}}", shift left, from=2-2, to=2-1]
    \end{tikzcd}\]
    By Lemma~\ref{lem:kleislirecognition}, $I$ is fully faithful, and thus $R_{\opname{KL}}$ is conservative. Since we are working inside $\sfSt_k$, $R_{\opname{KL}}$ is automatically colimit-preserving. Thus it is monadic, and agrees with the Eilenberg--Moore adjunction as desired.
\end{proof}

\begin{eg}
    It is useful to consider geometric examples of this. Let $Y$ be a smooth variety equipped with a map $f:Y\to \bbA^1/\bbG_m$. Write $D=Y\times_{\bbA^1/\bbG_m}(B\bbG_m)$, the associated generalized effective Cartier divisor.
    
    On the level of small categories, this yields an adjunction
    \[i^*:\Coh(Y)\tofrom \Coh(D):i_*\]
    which coincides with the Eilenberg--Moore construction as $i_*$ is seen to be monadic. However, the stable essential image of $i^*$ is $\Perf(D)$, noting that since $Y$ is smooth, $\Perf(Y)\isomto \Coh(Y)$. So the Kleisli adjunction associated to the monad $i^*i_*$ is
    \[i^*:\Perf(Y)\tofrom \Perf(D):i_*.\]
    So if $D$ is not smooth (for example take $Y=\pt$ and $f=0$) then the two adjunctions differ.

    On the other hand, in the large setting, both constructions output the adjunction,
    \[i^*:\QCoh(Y)\tofrom \QCoh(D):i_*.\]
    So this exhibits Proposition~\ref{prop:KLisEM} for the object $(\QCoh(Y),M_D := i_*i^*)$. An interesting consequence is that any other adjunction lifting this monad admits a universal map both to and from this one. Take for example the adjunction,
    \[i^*:\QCoh(Y)\tofrom \IndCoh(D):i_*,\]
    which yields the same monad $M_D$ on $\QCoh(Y)$. Then there are maps,
    \[\QCoh(D)\to \IndCoh(D)\to \QCoh(D).\]
    These are the adjoint functors $\Xi_D\adjto \Psi_D$ of \cite{gaitsgoryIndcoherentSheaves2012}.
\end{eg}

\subsection{Families of \texorpdfstring{$\two$}{2}-categories}\label{subsection:twocatfamily}
In this subsection, we prove some category theory results used in Section~\ref{section:families} and in Section~\ref{section:radon}.

\begin{defn}
    A local system of $\two$-categories on a space $X$ is a functor $X\to \wCat_{\two}$, where $X$ is viewed as a $\one$-category. A map between two local systems is a natural transformation of functors. We write $\Fun(X,\wCat_{\two})$ for the category of functors with morphisms given by natural transformations.
\end{defn}

If $X$ is connected and we fix a basepoint $x\in X$, then there is an identification,
\[X\simeq B\Omega_xX,\]
where $\Omega_xX$ be the loop group based at $x\in X$. In this way, a local system of $\two$-categories on $X$ is the same thing as a $\two$-category $\sfC_x$ with an action of $\Omega_xX$.

Consider the map of spaces, $\pi:X\to \pt$. Then, there is a pair of adjoint functors,
\[\pi^*:\wCat_{\two}\tofrom \Fun(X,\wCat_{\two}):\pi_*.\]
We will write $\pi_* := \Gamma(X,-)$. Under the above identification, this is the functor of $\Omega_xX$-fixed points.

\begin{prop}\label{prop:faithful}
    Let $X$ be a connected space and suppose we are given local systems of 2-categories, $\ul{\sfC},\ul{\sfD}$, together with a map $\opname{F}:\ul{\sfC}\to\ul{\sfD}$ between them. Fix $x\in X$, and suppose that $\opname{F}_x:\sfC_x\to\sfD_x$ is 2-faithful. Then, the natural functor
    \[\opname{I}:\Gamma(X,\ul{\sfC})\to \sfC_x\times_{\sfD_x}\Gamma(X,\ul{\sfD})\]
    is 2-fully faithful.

    Consider an object $\calE$ of the target given by $\calD\in \Gamma(X,\ul{\sfD})$ together with a lift of $\calD_x\in\sfD_x$ to some $\calC_x\in\sfC_x$. Then, $\calE$ lies in the essential image of $\opname{I}$ if and only if $\calC_x$ is a fixed point in the set of all lifts of $\calD_x$ along $\opname{F}_x$, under its natural $\pi_1(X,x)$-action.
\end{prop}

\begin{proof}
    As discussed above, we have an identification,
    \[\Gamma(X,\ul{\sfC})\simeq \sfC_x^{\Omega_xX}.\]
    We denote its objects by pairs $(\calC,\Xi)$ consisting of an object $\calC$ equipped with fixed point data $\Xi$. Given objects $(\calC_1,\Xi_1),(\calC_2,\Xi_2)$, the mapping category ${\sfC}(\calC_1,\calC_2)$ obtains a $\Omega_xX$-action, and
    \[{\sfC^{\Omega_xX}}((\calC_1,\Xi_1),(\calC_2,\Xi_2))\simeq {\sfC}(\calC_1,\calC_2)^{\Omega_xX}.\]
    We have similarly for $\Gamma(X,\ul{\sfD})\simeq \sfD^{\Omega_xX}$. We may organize our categories in the following diagram.
    % https://q.uiver.app/#q=WzAsNSxbMCwwLCJcXHNmQ194XntcXE9tZWdhX3hYfSJdLFsxLDAsIlxcc2ZFIl0sWzIsMCwiXFxzZkRfeF57XFxPbWVnYV94WH0iXSxbMSwxLCJcXHNmQ194Il0sWzIsMSwiXFxzZkRfeCJdLFswLDEsIlxcb3BuYW1le0l9Il0sWzEsMl0sWzEsM10sWzMsNCwiXFxvcG5hbWV7Rn1feCJdLFsyLDRdLFswLDMsIiIsMCx7ImN1cnZlIjoyfV1d
    \[\begin{tikzcd}
        {\sfC_x^{\Omega_xX}} & \sfE\pullback & {\sfD_x^{\Omega_xX}} \\
        & {\sfC_x} & {\sfD_x}
        \arrow["{\opname{I}}", from=1-1, to=1-2]
        \arrow[curve={height=12pt}, from=1-1, to=2-2]
        \arrow[from=1-2, to=1-3]
        \arrow[from=1-2, to=2-2]
        \arrow[from=1-3, to=2-3]
        \arrow["{\opname{F}_x}", from=2-2, to=2-3]
    \end{tikzcd}\]

    Now let $\calX_1=(\calC_1,\Xi_1),\calX_2=(\calC_2,\Xi_2)\in \sfC^{\Omega_xX}$. Then, we get the following diagram of hom categories.
    % https://q.uiver.app/#q=WzAsNSxbMCwwLCJcXEhvbV97XFxzZkNfeF57XFxPbWVnYV94WH19KFxcY2FsWF8xLFxcY2FsWF8yKVxcc2ltZXFcXEhvbV97XFxzZkNfeH0oXFxjYWxDXzEsXFxjYWxDXzIpXntcXE9tZWdhX3hYfSJdLFsxLDAsIlxcSG9tX3tcXHNmRX0oXFxvcG5hbWV7SX0oXFxjYWxYXzEpLFxcb3BuYW1le0l9KFxcY2FsWF8yKSkiXSxbMiwwLCJcXEhvbV97XFxzZkRfeH0oXFxvcG5hbWV7Rn0oXFxjYWxDXzEpLFxcb3BuYW1le0Z9KFxcY2FsQ18yKSlee1xcT21lZ2FfeFh9Il0sWzEsMSwiXFxIb21fe1xcc2ZDX3h9KFxcY2FsQ18xLFxcY2FsQ18yKSJdLFsyLDEsIlxcSG9tX3tcXHNmRF94fShcXG9wbmFtZXtGfShcXGNhbENfMSksXFxvcG5hbWV7Rn0oXFxjYWxDXzIpKSJdLFswLDEsIlxcb3BuYW1le0l9Il0sWzEsMl0sWzEsM10sWzMsNCwiXFxvcG5hbWV7Rn1feCJdLFsyLDRdLFswLDMsIiIsMCx7ImN1cnZlIjoyfV1d
    \[\begin{tikzcd}
        {{\sfC_x^{\Omega_xX}}(\calX_1,\calX_2)\simeq{\sfC_x}(\calC_1,\calC_2)^{\Omega_xX}} & {{\sfE}(\opname{I}(\calX_1),\opname{I}(\calX_2))}\pullback & {{\sfD_x}(\opname{F}(\calC_1),\opname{F}(\calC_2))^{\Omega_xX}} \\
        & {{\sfC_x}(\calC_1,\calC_2)} & {{\sfD_x}(\opname{F}(\calC_1),\opname{F}(\calC_2))}
        \arrow["{I}", from=1-1, to=1-2]
        \arrow[curve={height=12pt}, from=1-1, to=2-2]
        \arrow[from=1-2, to=1-3]
        \arrow[from=1-2, to=2-2]
        \arrow[from=1-3, to=2-3]
        \arrow["", from=2-2, to=2-3]
    \end{tikzcd}\]
    By assumption, the bottom horizontal arrow is fully faithful. By an analogous $\one$-categorical statement, $I$ is an equivalence of categories, so we deduce the 2-fully faithfulness.

    We now calculate the objects in the essential image of the functor $\opname{I}$. Let $\calD\in \Gamma(X,\ul{\sfD})$. Since $\opname{F}_x$ is faithful, the $(\infty,2)$-category of lifts of $\calD_x$ along $\opname{F}_x$ is in fact a partially ordered set (viewed as an $(\infty,2)$-category). Thus the action of the group $\Omega_xX$ factors through $\Omega_xX\to \pi_1(X,x)$. A lift $\calC_x$ of $\calD_x$ defines an object $\calE$ of $\sfC_x\times_{\sfD_x}\Gamma(X,\ul{\sfD})$.
    
    Suppose that $\calC_x$ is a fixed point for the $\pi_1(X,x)$-action. Then, for each $y\in X$, we can choose a path $\gamma$ from $x$ to $y$ in $X$, and obtain a lift $\calC_y^\gamma$ of $\calD_y$. The fixed point condition ensures that this is independent of $\gamma$, and so we obtain a preimage $\calC\in\sfC$. The converse is clear.
\end{proof}

\begin{cor}\label{cor:surjectivepi1}
    In the setting of Proposition~\ref{prop:faithful}, suppose that $U\to X$ is a map of connected spaces such that $x\in U$ and the induced map $\pi_1(U,x)\to \pi_1(X,x)$ is surjective. Then, the natural functor
    \[\opname{P}:\Gamma(X,\ul{\sfC})\to \Gamma(U,\ul{\sfC})\times_{\Gamma(U,\ul{\sfD})}\Gamma(X,\ul{\sfD})\]
    is an equivalence.
\end{cor}

\begin{proof}
    The essential image of the functor $\opname{I}$ of Proposition~\ref{prop:faithful} is given by fixed points for a $\pi_1(X,x)$-action, which is the same as being a fixed point under an associated $\pi_1(U,x)$-action. So the result follows.
\end{proof}

\begin{lem}\label{lem:faithfulpullback}
    Suppose we are given a pullback square of 2-categories as follows.
    \[\begin{tikzcd}
        {\sfC}\pullback & {\sfD} \\
        \sfC' & \sfD'
        \arrow["{\opname{F}}", from=1-1, to=1-2]
        \arrow["\opname{G}", from=1-1, to=2-1]
        \arrow["\opname{H}", from=1-2, to=2-2]
        \arrow["{\opname{F}'}", from=2-1, to=2-2]
    \end{tikzcd}\]
    If $\opname{F}'$ is 2-faithful, then so is $\opname{F}$.
\end{lem}

\begin{proof}
    Let $\calC_1,\calC_2\in\sfC$ be objects. Then, we have a pullback square of hom categories,
    \[\begin{tikzcd}
        {\sfC(\calC_1,\calC_2)}\pullback & {\sfD}(\opname{F}(\calC_1),\opname{F}(\calC_2)) \\
        {{\sfC'}(\opname{G}(\calC_1),\opname{G}(\calC_2))} & {{\sfD'}(\opname{HF}(\calC_1),\opname{HF}(\calC_2)).}
        \arrow[from=1-1, to=1-2]
        \arrow[from=1-1, to=2-1]
        \arrow[from=1-2, to=2-2]
        \arrow[from=2-1, to=2-2]
    \end{tikzcd}\]
    Since the bottom horizontal functor is an equivalence of categories, we deduce that the top horizontal functor is an equivalence, thus $\opname{F}$ is 2-faithful.
\end{proof}